% apendiceB.tex - Transcripción del avance del usuario (Apéndice B)
\chapter{Apéndice B: Diseño de placas base}

Este Apéndice trata sobre el diseño de placas base sujetas a carga axial y momentos flexionantes.

Está organizado como sigue:

B.1. Diseño de placas base
B.2. Resistencia al aplastamiento del concreto
B.3. Resistencia en flexión de la placa base
B.4. Resistencia en tensión del sistema de anclaje
B.5. Resistencia en cortante del sistema de anclaje
B.6. Interacción cortante y tensión en el sistema de anclaje

\section*{B.1 Diseño de placas base}

En las placas base sujetas a compresión, cortante o flexión, o una combinación de esas acciones, se deben revisar los siguientes modos de falla:

- Aplastamiento del concreto que soporta la placa base
- Fluencia de la placa base por flexión
- Falla de la placa en cortante
- Falla por una combinación de tensión y cortante del anclaje

\subsection*{B.1.1 En compresión axial}

Cuando la placa base de una columna soporta solo compresión axial, sus dimensiones en planta deben ser suficientes para transmitir esa fuerza al concreto, sin exceder su resistencia (estado límite de aplastamiento del concreto), y su grueso el adecuado para evitar el estado límite de flujo plástico de la placa.

Para el diseño de una placa base sometida a carga axial de compresión (Figura B.1.1) se recomienda el procedimiento siguiente:

1) Del análisis de la estructura se obtiene la carga axial de diseño en la base de la columna, $P_u$.

2) Se calcula el área requerida de la placa base, $A_{pl\,req}$.

\begin{center}
\textbf{aquí va la imagen 1}
\end{center}

3) Cálculo de las dimensiones de la placa base.

4) Se determina el momento de diseño máximo, $M_u$, considerando una distribución uniforme de presiones.

5) Grueso requerido de la placa base, $t_p$.

\subsection*{B.1.2 En tensión axial}

Para el diseño de una placa base en tensión (Figura B.1.2) se recomienda el procedimiento siguiente:

1) Del análisis de la estructura se obtiene la fuerza axial máxima de tensión en la base de la columna, $T_u$.

2) Se determina el número de anclas, $n$, requerido para resistir la fuerza axial máxima de tensión.

3) Se determinan las dimensiones en planta de la placa base, su grueso y la soldadura adecuada para transferir la fuerza de tensión.

4) Se revisa la resistencia del concreto correspondiente a los diferentes modos de falla posibles (arrancamiento de las anclas, desprendimiento del cono de concreto en tensión o separación al borde).

\begin{center}
\textbf{aquí va la imagen 2}
\end{center}

\subsection*{B.1.3 En cortante}

Las placas base en cortante se diseñan de acuerdo con el procedimiento siguiente:

1) Se determina el cortante máximo de diseño en la base de la columna, $V_u$.

2) Se selecciona el mecanismo de transferencia de la fuerza cortante de la columna a la cimentación (por fricción, por aplastamiento de la placa base o por cortante directo en las anclas).

\begin{center}
\textbf{aquí va la imagen 3}
\end{center}

\subsection*{Mecanismos (resumen)}

- Por fricción entre la placa base y el mortero de relleno o la superficie de concreto.
- Por aplastamiento de la placa base y parte de la columna, o de las llaves de cortante en el concreto de la cimentación.

\begin{center}
\textbf{aquí va la imagen 4}
\end{center}

\section*{B.2 a B.6 (resumen)}

Las secciones B.2 a B.6 contienen procedimientos de verificación de resistencia al aplastamiento del concreto, resistencia en flexión de la placa base, resistencia en tensión y cortante de anclajes, y la interacción entre cortante y tensión en los anclajes. En caso de diseño detallado se deben aplicar las fórmulas y criterios normativos correspondientes (ecuaciones y factores FR), tal como se indica en las secciones originales.

% Nota: las ecuaciones detalladas y tablas se han omitido en esta transcripción resumida; mantener las referencias normativas en `bibliografia/references.bib`.
